\documentclass[11pt]{article}
\usepackage[a4paper,margin=1in]{geometry}
\usepackage{amsmath,amssymb,amsthm,mathtools}
\usepackage[T1]{fontenc}
\usepackage{lmodern}
\usepackage{microtype}
\usepackage{enumitem}
\usepackage[hidelinks]{hyperref}

\newtheorem{proposition}{Proposition}[section]
\newtheorem{lemma}[proposition]{Lemma}
\newtheorem{definition}[proposition]{Definition}
\newtheorem{remark}[proposition]{Remark}

\newcommand{\dd}{\mathrm{d}}
\newcommand{\ii}{\mathrm{i}}
\newcommand{\half}{\frac12}
\newcommand{\Z}{\mathbb Z}
\newcommand{\C}{\mathbb C}
\newcommand{\R}{\mathbb R}
\newcommand{\cH}{\mathcal H}
\newcommand{\cI}{\mathcal I}
\newcommand{\cN}{\mathcal N}
\newcommand{\cS}{\mathcal S}
\newcommand{\e}{\mathrm e}

\numberwithin{equation}{section}

\begin{document}

\begin{center}
{\Large The Teukolsky--Starobinsky identities}\\[3pt]
\end{center}

\section{Kerr-AdS black holes}

\subsection{The radial case}

\subsubsection{The radial ODE and its solutions}

Fix parameters subextremal Kerr-AdS parameters $(a,M,L)$ with $M>0$, $L>0$. Define $r_+$ to be the largest root of 
\begin{equation} \label{eq:Delta-KAdS}
\Delta:=(r^2+a^2)\left(1+\frac{r^2}{L^2}\right)-2Mr.
\end{equation}
We also set
\begin{equation} \label{eq:Xi-KAdS}
\Xi:=1-\frac{a^2}{L^2}.
\end{equation}
For $\omega\in \mathbb C$ and $m\in \frac12\mathbb Z$, it will be convenient to define
\begin{equation}\label{eq:xi-KAdS}
        \xi:=-\ii\frac{2Mr_+}{r_+-r_-}(\omega-m\omega_+);\qquad 
        \omega_+:=\frac{a}{2Mr_+}.
\end{equation}
We say  $(\omega,m,\lambda)$ is an admissible frequency triple with respect to $a$ and $s$ if $(\omega,m,\lambda)\in \mathbb R^3$.

In this section, we discuss the radial ODE
\begin{equation}\label{eq:radial-ode-KAdS}
\begin{split}
&\Delta^{-s} \frac{d}{dr}\left(\Delta^{1+s}\frac{d}{dr}\right)R^{[s],(a\omega)}_{m\lambda}(r)+\frac{1}{\Delta}\left[\Xi^2(\omega(r^2+a^2)-am)^2 - is \Xi (\omega(r^2+a^2)-am) \frac{d\Delta}{dr}\right]R^{[s],(a\omega)}_{m\lambda}(r)\\
&\quad+\left[4is\omega \Xi r+\frac{2}{L^2}(1+3s+2s^2)r^2+s\left(1+\frac{a^2}{L^2}\right)-\lambda+2\Xi^2 am\omega-a^2\Xi^2\omega^2\right]R^{[s],(a\omega)}_{m\lambda}(r)=0\,,
\end{split}
\end{equation}
where $r\in (r_+,\infty)$.

The radial ODE \eqref{eq:radial-ode-KAdS} has regular singularities at every one of the four roots of $\Delta$, and another regular singularity at $r=\infty$. By standard asymptotic analysis, we can span the solution space by two linearly independent asymptotic solutions at each of the singularities $r=r_+,\infty$. Our basis is given by the following:
\begin{definition}\label{def:radial-basis-KAdS}
Fix $M>0$, $|a|\leq M$, $s\in \frac12\mathbb Z$ and an admissible frequency triple $(\omega,m,\lambda)$ with respect to $a$ and $s$.
\begin{enumerate}[label=\arabic*.]
\item Define $R^{[s]}_{\cH+}$ and $R^{[s]}_{\cH-}$ to be the unique classical solutions to the homogeneous radial ODE \eqref{eq:radial-ode-KAdS} as follows.
\begin{enumerate}[label=(\alph*)]
\item If $\omega-m\omega_+\neq 0$,
\begin{enumerate}[label=\roman*.]
\item $R^{[s]}_{\cH+}(r)(r-r_+)^{-\xi+s}$ and $R^{[s]}_{\cH-}(r)(r-r_+)^{\xi}$ are both smooth at $r=r_+$,
\item $\left((r-r_+)^{-\xi+s}R^{[s]}_{\cH+}\right)\big|_{r=r_+}=1$, and
 $\left((r-r_+)^{\xi}R^{[s]}_{\cH-}\right)\big|_{r=r_+}=1$;
\end{enumerate}
\item If $\omega=m\omega_+$, and $s\in \frac12\mathbb Z\backslash \mathbb Z_{> 0}$,
\begin{enumerate}[label=\roman*.]
\item $R^{[s]}_{\cH+}(r)(r-r_+)^{s}$ is smooth at $r=r_+$,
\item $\left((r-r_+)^{s}R^{[s]}_{\cH+}\right)\big|_{r=r_+}=1$,
\end{enumerate}
\item If $\omega=m\omega_+$, and $s\in \mathbb Z_{>0}$,
\begin{enumerate}[label=\roman*.]
\item $R^{[s]}_{\cH+}(r)$ is smooth at $r=r_+$,
\item $\left(R^{[s]}_{\cH+}\right)\big|_{r=r_+}=1$,
\end{enumerate}
\end{enumerate}
\item Define $R^{[s]}_{\cI,1}$ and $R^{[s]}_{\cI,2}$ to be the unique classical solutions to the homogeneous radial ODE \eqref{eq:radial-ode-KAdS} with boundary conditions
\begin{enumerate}[label=(\alph*)]
\item $R^{[s]}_{\cI,1}r^{1+2s}$ and  $R^{[s]}_{\cI,2}r^{2+2s}$ are smooth as functions of $r^{-1}$ as $r\to\infty$,
\item $\left(R^{[s]}_{\cI,1}r^{1+2s}\right)\big|_{r=\infty}=1$, and $\left(R^{[s]}_{\cI,2}r^{2+2s}\right)\big|_{r=\infty}=1$,
\item $r^{1+2s}R^{[s]}_{\cI,1}-1=O(r^{-2})$ and $r^{2+2s}R^{[s]}_{\cI,2}-1=O(r^{-2})$ as $r\to \infty$.
\end{enumerate}
\end{enumerate}
\end{definition}

\begin{remark} The condition $r^{2+2s}R^{[s]}_{\cI,2}-1=O(r^{-2})$ follows from the previous constraints on $R^{[s]}_{\cI,2}$, through the recurrence relation for the Frobenius representation of this solution. On the contrary, the condition $r^{1+2s}R^{[s]}_{\cI,1}-1=O(r^{-2})$ is a {\normalfont choice} that is made so that $R^{[s]}_{\cI,1}$ is uniquely specified.
\end{remark}


\subsubsection{The radial Teukolsky--Starobinsky identities}

In this section, we will introduce the Teukolsky--Starobinsky identities. We begin by defining the differential operators
\begin{equation}\label{eq:Dpm-KAdS}
        D^{\pm}_n=\frac{\dd}{\dd r}\pm\ii\left(\frac{\Xi\omega(r^2+a^2)}{\Delta}-\frac{a\Xi m}{\Delta}\right)+\frac{n}{\Delta}\frac{d\Delta}{dr}.
\end{equation}

The radial Teukolsky--Starobinsky identities are differential identities that relate solutions of the  radial ODE \eqref{eq:radial-ode} with spin $+s$ and spin $-s$. We will present a rigorous statement:

\begin{proposition}[Radial Teukolsky--Starobinsky identities]\label{prop:radial-ts-KAdS}
Fix subextremal $(a,M,L)$ Kerr-AdS parameters and $s\in \frac12\mathbb Z_{\geq0}$. Let $(\omega,m,\lambda)$ be an admissible frequency triple with respect to $s$ and $a$. There are constants $\mathfrak C^{[\pm s]}_{ji}\in\mathbb C$, for $i,j\in\{1,2\}$, which depend only on $s$, $(a,M,L)$ and $(\omega,m,\lambda)$, such that 
\begin{align}
\begin{split}
\Delta^s(D^+ _0)^{2s}\left(\Delta^s R^{[+ s]}_{\cI,j}\right)
&=\mathfrak C^{[+ s]}_{j1}R^{[- s]}_{\cI,1}
  +\mathfrak C^{[+ s]}_{j2}R^{[- s]}_{\cI,2}\\
\Delta^s(D^- _0)^{2s}\left( R^{[-s]}_{\cI,j}\right)
&=\mathfrak C^{[- s]}_{j1}\Delta^sR^{[+ s]}_{\cI,1}
  +\mathfrak C^{[- s]}_{j2}\Delta^sR^{[+ s]}_{\cI,2}
\\
  \end{split}
\label{eq:TS-identities-KAdS}
\end{align}
for $j=1,2$. Furthermore, we have \texttt{at least for $s\leq 5/2$ [to check other cases]}
\begin{align*}
\left[\mathfrak C^{[0]}]\right]&=\begin{pmatrix}
1&0\\
0&1\\
\end{pmatrix},\\
\left[\mathfrak C^{[\pm s]}\right]&=(\pm 1)^{\sin^2(\pi s)}L^{\mp s}
\begin{pmatrix}
  \ii^{\sin^2(\pi s)}(\mathfrak p_{s,0}\mp \ii \mathfrak q_s) & \pm \ii^{\cos^2(\pi s)} L \mathfrak p_{s,2}\\
\pm  \ii^{\cos^2(\pi s)}L^{-1}\mathfrak p_{s,1}& \ii^{\sin^2(\pi s)}(\mathfrak p_{s,0}\pm \ii \mathfrak q_s) 
\end{pmatrix}, \qquad s\neq 0.
\end{align*}
%i.e.\
%\begin{align*}
%\left[\mathfrak C^{[0]}]\right]&=\begin{pmatrix}
%1&0\\
%0&1\\
%\end{pmatrix},\\
%\left[\mathfrak C^{[\pm \frac12]}\right]&= \mp L^{\mp 1}\begin{pmatrix}
% i \mathfrak p_{\frac12,0} & \pm L \mathfrak p_{\frac12,2}\\
%\pm L^{-1}\mathfrak p_{\frac12,1}&  i\mathfrak p_{\frac12,0}
%\end{pmatrix},\\
%\left[\mathfrak C^{[\pm 1]}\right]&= L^{\mp 2}\begin{pmatrix}
%\mathfrak p_{1,0} & \pm  i L\mathfrak p_{1,2}\\
%\pm iL^{-1}\mathfrak p_{1,1}& \mathfrak p_{1,0}
%\end{pmatrix},\\
%%
%\left[\mathfrak C^{[\pm \frac32]}\right]&=\mp  L^{\mp 3}\begin{pmatrix}
%i\mathfrak p_{\frac32,0}\pm \mathfrak q_{\frac32} &\pm L\mathfrak p_{\frac32,2}\\
%\pm  L^{-1}\mathfrak p_{\frac32,1}& i\mathfrak p_{\frac32,0}\mp \mathfrak q_{\frac32}
%\end{pmatrix},\\
%%
%\left[\mathfrak C^{[\pm 2]}\right]
%&=
%L^{\mp 4}
%\begin{pmatrix}
%\mathfrak p_{2,0}\mp i\mathfrak q_2
%&
%\pm iL\mathfrak p_{2,2}
%\\
%\pm iL^{-1}\mathfrak p_{2,1} 
%&
%\mathfrak p_{2,0}\pm i\mathfrak q_2
%\end{pmatrix},\\
%%
%\left[\mathfrak C^{[\pm \frac52]}\right]&=\mp  L^{\mp 5}\begin{pmatrix}
%i\mathfrak  p_{\frac52,0}\pm \mathfrak q_{\frac52} &\pm  L\mathfrak p_{\frac52,2}\\
%\pm L^{-1}\mathfrak p_{\frac52,1}& i\mathfrak p_{\frac52,0}\mp \mathfrak q_{\frac52}
%\end{pmatrix}.
%\end{align*}
Here, $\hat \Lambda=\lambda+a^2\Xi^2\omega^2-2a\Xi^2m\omega+|s|\left(1+\frac{a^2}{L^2}\right)$, and  $\mathfrak q_{s},\mathfrak p_{s,0},\mathfrak p_{s,1},\mathfrak p_{s,2}$ are given by
\begin{itemize}
\item for $s=1/2$,
\begin{align*}
\mathfrak q_{\frac12}=0, \qquad\mathfrak{p}_{\frac12,0}= -L\Xi\omega,\qquad 
\mathfrak{p}_{\frac12,1}= 1,\qquad
\mathfrak{p}_{\frac12,2}= \hat\Lambda-L^2\Xi^2\omega^2,
\end{align*}
with $\mathfrak{p}_{\frac12,0}^2+\mathfrak{p}_{\frac12,1}\mathfrak{p}_{\frac12,2}=\hat \Lambda$;
\item for $s=1$,
\begin{gather*}
\mathfrak q_{1}=0, \qquad\mathfrak{p}_{1,0}= \hat\Lambda-2L^2\Xi^2\omega^2,\qquad
\mathfrak{p}_{1,1}= -2L\Xi\omega,\\
\mathfrak{p}_{1,2}= 2L^3\Xi^3\omega^3-2L\hat\Lambda\Xi\omega+\frac{2a}{L}(a\Xi\omega-\Xi m),
\end{gather*}
with $\mathfrak{p}_{1,0}^2+\mathfrak{p}_{1,1}\mathfrak{p}_{1,2}=\hat \Lambda^2+4a\Xi^2 m\omega-4a^2\Xi^2\omega^2$;
\item for $s=3/2$,
\begin{align*}
\mathfrak q_{\frac32}&=\frac{2M}{L}, \\
\mathfrak p_{\frac32,0}
&= -\frac{4a}{L}\Xi m + 3\frac{a^2}{L}\Xi\omega
+ L\Xi\omega(-1 - 3\hat\Lambda + 4L^2\Xi^2\omega^2),\\
\mathfrak p_{\frac32,1}
&=  1 + \hat\Lambda - 4L^2\Xi^2\omega^2+\frac{a^2}{L^2},\\
\mathfrak p_{\frac32,2}
&= \hat\Lambda^2 - 8a\Xi^2m\omega
- L^2(1 + 5\hat\Lambda)\Xi^2\omega^2
+ 4L^4\Xi^4\omega^4
- \frac{a^2}{L^2}(4 - 7L^2\Xi^2\omega^2),
\end{align*}
with
\begin{align*} \mathfrak p_{\frac32,0}^2+\mathfrak p_{\frac32,1}\mathfrak p_{\frac32,2}&=\hat\Lambda^2\left(\hat\Lambda+1+\frac{a^2}{L^2}\right)-4\hat\Lambda\left(4a^2\Xi^2\omega^2-4a\Xi^2m\omega+\frac{a^2}{L^2}\right)\\
&+16a^2\Xi^2\omega^2\left(1+\frac{a^2}{L^2}\right) -\frac{4a^2}{L^2}\left(1+\frac{a^2}{L^2}-4\Xi^2m^2+8a\Xi^2m\omega\right);
\end{align*}
\item for $s=2$,
\begin{align*}
\mathfrak q_{2}&=12M\Xi\omega, \\
\mathfrak p_{2,0}
&= 2\hat\Lambda + \hat\Lambda^2 - 20a\Xi^2m\omega
- 8L^2(1+\hat\Lambda)\Xi^2\omega^2
+ 8L^4\Xi^4\omega^4
+ 2a^2\left(\frac{\hat\Lambda - 6}{L^2} + 6\Xi^2\omega^2\right),\\
\mathfrak p_{2,1}
&= -8\frac{a}{L}\Xi m - 4L(2+\hat\Lambda)\Xi\omega
+ 8L^3\Xi^3\omega^3,\\
\mathfrak p_{2,2}
&= -4\hat\Lambda(1+\hat\Lambda)L\Xi\omega
+ 4L^3(2+3\hat\Lambda)\Xi^3\omega^3
- 8L^5\Xi^5\omega^5
- \frac{12a}{L}\hat\Lambda\Xi m
+ 32aL\Xi^3m\omega^2
\\
&\qquad + \frac{8a^2}{L}(3+\hat\Lambda)\Xi\omega
- 24a^2L\Xi^3\omega^3,
\end{align*}
with
\begin{align*}
\mathfrak{p}_{2,0}^2+\mathfrak{p}_{2,1}\mathfrak{p}_{2,2}&=\hat\Lambda^2\left(\hat\Lambda+2\left(1+\frac{a^2}{L^2}\right)\right)^2
-8\hat\Lambda^2\left(5a^2\Xi^2\omega^2-5a\Xi^2m\omega+3\frac{a^2}{L^2}\right)\\
&\qquad
+48\hat\Lambda\left(
\left(1+\frac{a^2}{L^2}\right)a^2\Xi^2\omega^2
+a\Xi^2m\omega\left(1-3\frac{a^2}{L^2}\right)
+2\frac{a^2}{L^2}\Xi^2 m^2
\right)\\
&\qquad
-48\frac{a^2}{L^2}\hat\Lambda \left(1+\frac{a^2}{L^2}\right)
+144\left(a\Xi^2m\omega-a^2\Xi^2\omega^2+\frac{a^2}{L^2}\right)^2;
\end{align*}
\item for $s=5/2$,
\begin{align*}
\mathfrak q_{\frac52}
&= \frac{6M}{L^3}\left(6a^2 + L^2(6+\hat\Lambda - 8L^2\Xi^2\omega^2)\right),\\
\mathfrak p_{\frac52,0}
&= -24\frac{a^3}{L^3}\Xi m + 15\frac{a^4}{L^3}\Xi\omega
+ 4\frac{a}{L}\Xi m(-6 - 5\hat\Lambda + 20L^2\Xi^2\omega^2)\\
&\qquad
+ \frac{a^2}{L}(90\Xi\omega - 40L^2\Xi^3\omega^3)\\
&\qquad
- L\Xi\omega\left(
9 + 5\hat\Lambda^2 - 40L^2\Xi^2\omega^2
+ 16L^4\Xi^4\omega^4
- 20\hat\Lambda(-1 + L^2\Xi^2\omega^2)
\right),\\
\mathfrak p_{\frac52,1}
&= \frac{9a^4}{L^4} - 48a\Xi^2m\omega
+ 2\frac{a^2}{L^2}(-9 + 3\hat\Lambda + 4L^2\Xi^2\omega^2)\\
&\qquad
+\left(
9 + \hat\Lambda^2 - 40L^2\Xi^2\omega^2
+ 16L^4\Xi^4\omega^4
+ \hat\Lambda(6 - 12L^2\Xi^2\omega^2)
\right),\\
\mathfrak p_{\frac52,2}
&= - \frac{144M^2}{L^2} + \hat\Lambda^2(4+\hat\Lambda) 
+ 80\frac{a^3}{L^2}\Xi^2m\omega - 25\frac{a^4}{L^2}\Xi^2\omega^2\\
&\qquad
- L^2(9 + 34\hat\Lambda + 13\hat\Lambda^2)\Xi^2\omega^2
+ 4L^4(10 + 7\hat\Lambda)\Xi^4\omega^4
- 16L^6\Xi^6\omega^6\\
&\qquad
+ 8a\Xi^2m\omega(-6 - 9\hat\Lambda + 14L^2\Xi^2\omega^2)\\
&\qquad
+ \frac{a^2}{L^2}\left(
4\hat\Lambda^2 - 64\Xi^2m^2
+ 162L^2\Xi^2\omega^2 - 72L^4\Xi^4\omega^4
- \hat\Lambda(48 - 38L^2\Xi^2\omega^2)
\right).
\end{align*}
with $\mathfrak p_{\frac52,0}^2+\mathfrak p_{\frac52,1}\mathfrak p_{\frac52,2}$ \texttt{\textbf{NOT} equal to the angular TS constant and hence possibly not positive [?]}.
\end{itemize}



\end{proposition}

In order to prove Proposition \ref{prop:radial-ts-KAdS}, we quote the following two results, obtained from adapting [KMW89] to the asymptotically AdS case:

\begin{lemma}\label{lem:P-ode-KAdS} For solutions of the homogeneous radial ODE~\eqref{eq:radial-ode-KAdS} with spin $\pm s$, we set
\begin{align*}
P^{[+ s]}:=\Delta^s R^{[+s]}, \quad P^{[-s]}:=R^{[-s]}.
\end{align*}
Then, we have
\begin{equation}\label{eq:P-ode-KAdS}
\begin{split}
&\left[\Delta D^{\mp}_{1-s}D^{\pm}_0\pm2(2s-1)\ii\Xi\omega r+\frac{2r^2}{L^2}(1-3|s|+2s^2)\right]P^{[\pm s]}(r)\\
&\quad=\left(\lambda+|s|\left(1+\frac{a^2}{L^2}\right)+a^2\Xi^2\omega^2-2a\Xi^2m\omega\right)P^{[\pm s]}(r).
\end{split}
\end{equation}
\end{lemma}

By induction on $2s$, one can also show that:

\begin{lemma}\label{lem:radial-commutation-KAdS}
The operators $D^{\pm}_0$, defined in \eqref{eq:Dpm-KAdS}, satisfy the commutation relation
\begin{align}\label{eq:radial-commutation-KAdS}
&\Delta^s(D^{\mp}_0)^{2s}
\left[\Delta D^{\pm}_{1-s}D^{\mp}_0\mp2(2s-1)\ii\Xi\omega r+\frac{2r^2}{L^2}(1-3|s|+2s^2)\right]
\notag\\
&\hspace{35mm}=
\left[\Delta D^{\mp}_{1-s}D^{\pm}_0\pm2(2s-1)\ii\Xi\omega r+\frac{2r^2}{L^2}(1-3|s|+2s^2)\right]
\Delta^s(D^{\mp}_0)^{2s}.
\end{align}
\end{lemma}

Combining the two lemmas, we see that, if $P^{[\pm s]}$ solves \eqref{eq:P-ode} with spin $\pm$, then $\Delta^s(D^{\pm}_0)^{2s}P^{[\pm s]}$ is a solution of the homogeneous radial ODE \eqref{eq:P-ode} for spin $\mp s$.

\begin{proof}[Proof of Proposition \ref{prop:radial-ts}]
The statement is trivial for $s=0$. For $s\neq0$, we obtain the result by considering  the Frobenius representations of $R^{[\pm s]}_{\cI,1}$ and $R^{[\pm s]}_{\cI,2}$, on which we can act directly with the Teukolsky--Starobinsky operators. In terms of $x=1/r$, these are
\begin{align*}
R^{[\pm s]}_{\cI,1}=x^{1\pm 2s}\left(1+\sum_{k=2}^\infty c_{k,1}^{[\pm s]}x^k\right),\\
R^{[\pm s]}_{\cI,2}=x^{2\pm 2s}\left(1+\sum_{k=2}^\infty c_{k,2}^{[\pm s]}x^k\right).
\end{align*}
[COMPLETE ME]
\end{proof}

\subsubsection{An energy current}

In this section, we derive an energy current for the radial ODE~\eqref{eq:radial-ode-KAdS}. To do, it is convenient to recast it in terms of another variable: writing $u^{[s],(a\omega)}_{m\lambda}= (r^2+a^2)^{1/2}\Delta^{s/2}R^{[s],(a\omega)}_{m\lambda}$, we have
\begin{gather}
    \partial_{r^*}^2 u^{[s],(a\omega)}_{m\lambda} + (\Xi^2\omega^2-V^{[s],(a\omega)}_{m\lambda})u^{[s],(a\omega)}_{m\lambda}=0,  \label{eq:radial-ode-u-KAdS}\\
    \begin{aligned}
    V^{[s],(a\omega)}_{m\lambda}&= \frac{2am\Xi^2\omega}{r^2+a^2}
-\frac{a^2\Xi^2m^2}{(r^2+a^2)^2}+is\frac{\partial_r\Delta}{r^2+a^2}\left(
\Xi\omega
-\frac{a\Xi m}{(r^2+a^2)}
\right)\\
&\quad-\frac{\Delta}{(r^2+a^2)^2}\left(
\frac{2 r^2}{L^2}
-\lambda+2\Xi^2am\omega-a^2\Xi^2\omega^2
+4is\Xi\omega r\right)
+\partial_{r^*}\left(\frac{r\Delta}{(r^2+a^2)^2}\right)\\
&\quad
+s^2\left[
\frac{(\partial_r\Delta)^2}{4(r^2+a^2)^2}
-\frac{4r^2\Delta}{L^2(r^2+a^2)^2}
\right] +\left(\frac{r\Delta}{(r^2+a^2)^2}\right)^2,
\end{aligned}\notag
\end{gather}
where we have introduced a coordinate $r^*\colon (-\infty, \pi/2)$ defined via
\begin{align*}
    \frac{dr^*}{dr}=\frac{r^2+a^2}{\Delta}, \quad r^*(r)\to \pi/2 \text{~as $r\to \infty$.}
\end{align*}

Our energy current is based on the following conservation law:

\begin{lemma}\label{lem:Wronskian-conservation-KAdS} Fix $(a,M,L)$ subextremal and $s\in\frac12\mathbb Z_{\geq 0}$. If $(\omega,m,\lambda)$ is an admissible frequency triple with respect to $s$ and $a$, then
\begin{align*}
\frac{d}{dr^*}W[\overline{u^{[+s]}},u^{[-s]}]=0, \qquad W[\overline{u^{[+s]}},u^{[-s]}]= \overline{u^{[+s]}}\frac{d}{dr^*}u^{[-s]}- u^{[-s]}\frac{d}{dr^*}\overline{u^{[+s]}}
\end{align*}
\end{lemma}

We then have
\begin{proposition}[Energy current] 
[To FILL]
\end{proposition}

\begin{proof}
Note that there are constants $a_{\cI,1}^{[\pm s]}$, $a_{\cI,2}^{[\pm s]}$ such that 
\begin{align*}
u^{[\pm s]} = a_{\cI,1}^{[\pm s]}\Delta^{\pm s/2}(r^2+a^2)^{1/2}R_{\cI,1}^{[\pm s]}+ a_{\cI,2}^{[\pm s]}\Delta^{\pm s/2}(r^2+a^2)^{1/2}R_{\cI,2}^{[\pm s]};
\end{align*}
so we have 
\begin{align*}
\lim_{r\to \infty}W[\overline{u^{[+s]}},u^{[-s]}](r^*(r))= \frac{1}{L^2}(a_{\cI,1}^{[-s]},\overline{a_{\cI,2}^{[+s]}}-a_{\cI,2}^{[-s]},\overline{a_{\cI,1}^{[+s]}}).
\end{align*}

To obtain a conservation law for $\pm s$, we apply Lemma~\ref{lem:Wronskian-conservation-KAdS} with $u^{[\mp s]}$ obtained from $u^{[\pm s]}$ via the Teukolsky--Starobinsky map in Proposition~\ref{prop:radial-ts-KAdS}. In this case, $a_{\cI,1}^{[\pm s]}$ and $a_{\cI,2}^{[\pm s]}$ are free but  $a_{\cI,1}^{[\mp s]}$ and $a_{\cI,2}^{[\mp s]}$ are determined by the aforementioned proposition:
\begin{align*}
a_{\cI,1}^{[\mp s]}&=\mathfrak C_{11}^{[\pm s]}a_{\cI,1}^{[\pm s]}+\mathfrak C_{21}^{[\pm s]}a_{\cI,2}^{[\pm s]}\\
a_{\cI,2}^{[\mp s]}&=\mathfrak C_{12}^{[\pm s]}a_{\cI,1}^{[\pm s]}+\mathfrak C_{22}^{[\pm s]}a_{\cI,2}^{[\pm s]}
\end{align*}
Then, we have
\begin{align*}
\lim_{r\to \infty}W[\overline{u^{[+s]}},u^{[-s]}](r^*(r))
=
\frac{1}{L^2}\Big(
&-\mathfrak C_{12}^{[+s]}
\left|a_{\cI,1}^{[+s]}\right|^2
+\mathfrak C_{21}^{[+s]}
\left|a_{\cI,2}^{[+s]}\right|^2 \\
&+\mathfrak C_{11}^{[+s]}
a_{\cI,1}^{[+s]}\overline{a_{\cI,2}^{[+s]}}
-\mathfrak C_{22}^{[+s]}
a_{\cI,2}^{[+s]}\overline{a_{\cI,1}^{[+s]}}
\Big),
\end{align*}
and
\begin{align*}
\lim_{r\to \infty}
W[\overline{u^{[+s]}},u^{[-s]}](r^*(r))
=
\frac{1}{L^2}\Big(
&\overline{\mathfrak C_{12}^{[-s]}}
\left|a_{\cI,1}^{[-s]}\right|^2
-\overline{\mathfrak C_{21}^{[-s]}}
\left|a_{\cI,2}^{[-s]}\right|^2 \\
&+\overline{\mathfrak C_{22}^{[-s]}}
a_{\cI,1}^{[-s]}\overline{a_{\cI,2}^{[-s]}}
-\overline{\mathfrak C_{11}^{[-s]}}
a_{\cI,2}^{[-s]}\overline{a_{\cI,1}^{[-s]}}
\Big),
\end{align*}
respectively. For instance, using Proposition~\ref{prop:radial-ts-KAdS}, we have
%\begin{align*}
%& L^{-3}\left(
%L\mathfrak p_{\frac12,2}|a_{\cI,1}^{[+\frac12]}|^2
%- L^{-1}\mathfrak p_{\frac12,1}|a_{\cI,2}^{[+\frac12]}|^2
%+ 2\mathfrak p_{\frac12,0}\Im(a_{\cI,1}^{[+\frac12]}\overline{a_{\cI,2}^{[+\frac12]}})
%\right),
%\\
%& \ii L^{-4}\left(
%- L\mathfrak p_{1,2}|a_{\cI,1}^{[+1]}|^2
%+ L^{-1}\mathfrak p_{1,1}|a_{\cI,2}^{[+1]}|^2
%+ 2\mathfrak p_{1,0}\Im(a_{\cI,1}^{[+1]}\overline{a_{\cI,2}^{[+1]}})
%\right),
%\\
%& L^{-5}\left(
%L\mathfrak p_{\frac32,2}|a_{\cI,1}^{[+\frac32]}|^2
%- L^{-1}\mathfrak p_{\frac32,1}|a_{\cI,2}^{[+\frac32]}|^2
%+ 2\mathfrak p_{\frac32,0}\Im(a_{\cI,1}^{[+\frac32]}\overline{a_{\cI,2}^{[+\frac32]}})
%- \frac{4M}{L}\Re(a_{\cI,1}^{[+\frac32]}\overline{a_{\cI,2}^{[+\frac32]}})
%\right),
%\\
%&\ii L^{-6}\left(
%-L\mathfrak p_{2,2}|a_{\cI,1}^{[+2]}|^2
%+ L^{-1}\mathfrak p_{2,1}|a_{\cI,2}^{[+2]}|^2
%+ 2\mathfrak p_{2,0}\Im(a_{\cI,1}^{[+2]}\overline{a_{\cI,2}^{[+2]}})
%- 24M\Xi\omega\Re(a_{\cI,1}^{[+2]}\overline{a_{\cI,2}^{[+2]}})
%\right),
%\\
%& L^{-7}\left(
%L\mathfrak p_{\frac52,2}|a_{\cI,1}^{[+\frac52]}|^2
%- L^{-1}\mathfrak p_{\frac52,1}|a_{\cI,2}^{[+\frac52]}|^2
%+ 2\mathfrak p_{\frac52,0}\Im(a_{\cI,1}^{[+\frac52]}\overline{a_{\cI,2}^{[+\frac52]}})
%- 2\mathfrak q_{\frac52}\Re(a_{\cI,1}^{[+\frac52]}\overline{a_{\cI,2}^{[+\frac52]}})
%\right),
%\end{align*}
\begin{align*}
\lim_{r\to \infty}W[\overline{u^{[+s]}},u^{[-s]}](r^*(r))=\ii^{\cos^2(\pi s)}L^{-(2s+2)}&\left(
-(-1)^{2s}L\mathfrak p_{s,2}|a_{\cI,1}^{[+s]}|^2
+(-1)^{2s} L^{-1}\mathfrak p_{s,1}|a_{\cI,2}^{[+s]}|^2\right.
\\
&\qquad\left.+ 2\mathfrak p_{s,0}\Im(a_{\cI,1}^{[+s]}\overline{a_{\cI,2}^{[+s}})
- 2\mathfrak q_{s}\Re(a_{\cI,1}^{[+s]}\overline{a_{\cI,2}^{[+s]}})
\right)
\end{align*}
in the first case, and [COMPLETE ME WITH SECOND CASE]
\end{proof}

\begin{lemma} Suppose $\mathfrak P_s:=\mathfrak p_{s,0}^2+\mathfrak p_{s,1}\mathfrak p_{s,2}>0$. Then, 
\begin{align*}
\begin{pmatrix}
\mathfrak p_{s,0}\pm\sqrt{\mathfrak P_s}
&
\ii (-1)^{2s }L^{-1}\mathfrak p_{s,1}
\\
\ii (-1)^{2s} L\mathfrak p_{s,2}
&
\mathfrak p_{s,0}\mp\sqrt{\mathfrak P_s}
\end{pmatrix}\binom{a_{\cI,1}^{[+s]}}{a_{\cI,2}^{[+s]}}
=0 \implies \lim_{r\to \infty}W[\overline{u^{[+s]}},u^{[-s]}](r^*(r))=0.
\end{align*}
\texttt{The converse is not true.}

[COMPLETE ME WITH STATEMENT FOR NEGATIVE SPIN]
\end{lemma}

\begin{proof} A simple computation yields
\begin{align*}
\begin{pmatrix}
\mathfrak p_{s,0}\pm\sqrt{\mathfrak P_s}
&
\ii (-1)^{2s }L^{-1}\mathfrak p_{s,1}
\\
\ii (-1)^{2s} L\mathfrak p_{s,2}
&
\mathfrak p_{s,0}\mp\sqrt{\mathfrak P_s}
\end{pmatrix}\binom{a_{\cI,1}^{[+s]}}{a_{\cI,2}^{[+s]}}=0 \\
\implies 
\begin{dcases}
(-1)^{2s}L^{-1}\mathfrak p_{s,1}|a_{\cI,2}^{[+s]}|^2=i(\mathfrak p_{s,0}\pm\sqrt{\mathfrak P})a_{\cI,1}^{[+s]}\overline{a_{\cI,2}^{[+s]}}\\
(-1)^{2s}L^{-1}\mathfrak p_{s,2}|a_{\cI,1}^{[+s]}|^2=i(\mathfrak p_{s,0}\mp\sqrt{\mathfrak P})\overline{a_{\cI,1}^{[+s]}}a_{\cI,2}^{[+s]}
\end{dcases}.
\end{align*}
Subtracting the two equations, we get
\begin{align*}
-(-1)^{2s}L\mathfrak p_{s,2}|a_{\cI,1}^{[+s]}|^2
+(-1)^{2s} L^{-1}\mathfrak p_{s,1}|a_{\cI,2}^{[+s]}|^2=-2\mathfrak p_{s,0}\Im(a_{\cI,1}^{[+s]}\overline{a_{\cI,2}^{[+s}})+2\ii \sqrt{\mathfrak P_s} \Re(a_{\cI,1}^{[+s]}\overline{a_{\cI,2}^{[+s}});
\end{align*}
on the other hand, taking imaginary parts of the two equations, we obtain
\begin{align*}
(\mathfrak p_{s,0}\mp\sqrt{\mathfrak P_s})\Re(\overline{a_{\cI,1}^{[+s]}}a_{\cI,2}^{[+s]})=0=(\mathfrak p_{s,0}\pm\sqrt{\mathfrak P_s})\Re(\overline{a_{\cI,1}^{[+s]}}a_{\cI,2}^{[+s]})\implies \Re(\overline{a_{\cI,1}^{[+s]}}a_{\cI,2}^{[+s]})=0,
\end{align*}
since $\sqrt{\mathfrak P_s}\neq 0$. Thus, we have
\begin{align*}
\lim_{r\to \infty}W[\overline{u^{[+s]}},u^{[-s]}](r^*(r))=-2\ii^{\cos^2(\pi s)}L^{-(2s+2)}
\mathfrak q_{s}\Re(a_{\cI,1}^{[+s]}\overline{a_{\cI,2}^{[+s]}})=0.
\end{align*}
\end{proof}


\subsection{Comparison with Graf--Holzegel}

In comparing the section above with the paper of Graf--Holzegel, it is important to note that 
\begin{align*}
L&=1/k_{\text{GH}},\\
\hat \Lambda&=\lambda_{m\ell,\text{~GH}}^{\omega},\\
R_{m\lambda}^{[\pm 2],(a\omega)}&= \frac{(r^2+a^2)^{3/2}}{\Delta^{(1\pm 1)}}R^{m\ell}_{[\pm 2],\omega, \text{~GH}},
\end{align*}
where the subscript GH refers to the Graf--Holzegel variables.

The Robin boundary conditions of Graf--Holzegel's appendix are, in our notation,
\[
\mathcal R_\pm
\binom{a_{\cI,1}^{[+2]}}{a_{\cI,2}^{[+2]}}
=0, \quad \mathcal R_\pm
\binom{a_{\cI,1}^{[-2]}}{a_{\cI,2}^{[-2]}}
=0,
\]
where
\[
\mathcal R_\pm
:=
\begin{pmatrix}
\mathfrak p_{2,0}\pm\sqrt{\mathfrak p_{2,0}^2+\mathfrak p_{2,1}\mathfrak p_{2,2}}
&
\ii L^{-1}\mathfrak p_{2,1}
\\
\ii L\mathfrak p_{2,2}
&
\mathfrak p_{2,0}\mp\sqrt{\mathfrak p_{2,0}^2+\mathfrak p_{2,1}\mathfrak p_{2,2}}
\end{pmatrix}.
\]

\subsection{Comparison with Wang--Herdeiro--Jing/Sampaio}

In comparing the section above with the paper of  Wang--Herdeiro--Sampaio, it is important to note that (comparing the radial ODEs) 
\begin{align*}
\hat\Lambda&=\lambda_{\text{WHS}}-\frac{a^2}{L^2}\\
R_{m\lambda}^{[\pm 1],(a\omega)}&= R_{\pm 1, \text{~WHS}},
\end{align*}
where the subscript WHS refers to the Wang--Herdeiro--Sampaio variables.

The Robin boundary conditions of Wang--Herdeiro--Sampaio appendix are, in our notation, [FILL ME].

[ADD COMPARISON TO Wang--Herdeiro--Jing]

\section{Kerr black holes}



\subsection{The radial case}



\subsubsection{The radial ODE and its solutions}
Fix parameters $(a,M)$ with $M>0$ and $|a|\leq M$ and define
\begin{equation}\label{eq:rpm}
        r_{\pm}:=M\pm \sqrt{M^2-a^2}.
\end{equation}
We write
\begin{equation}\label{eq:Delta}
        \Delta:=r^2-2Mr+a^2=(r-r_+)(r-r_-).
\end{equation}
For $\omega\in \mathbb C$ and $m\in \frac12\mathbb Z$, it will be convenient to define
\begin{equation}\label{eq:xi-beta}
        \xi:=-\ii\frac{2Mr_+}{r_+-r_-}(\omega-m\omega_+), \text{~~if $|a|<M$};\qquad
        \beta:=2\ii M^2(\omega-m\omega_+);\qquad
        \omega_+:=\frac{a}{2Mr_+}.
\end{equation}
We say  $(\omega,m,\lambda)$ is an admissible frequency triple with respect to $a$ and $s$ if $(\omega,m,\lambda)\in \mathbb R^3$,  $\omega\neq 0$ and, for $|a|=M$, if $\omega\neq m\omega_+$ too.

In this section, we discuss the radial ODE
\begin{align}
&\Delta^{-s}\frac{\dd}{\dd r}\left(\Delta^{s+1}\frac{\dd}{\dd r}\right)
        R^{[s],(a\omega)}_{m\lambda}(r)
 +\left(\frac{[\omega(r^2+a^2)-am]^2-2\ii s(r-M)[\omega(r^2+a^2)-am]}{\Delta}\right)
        R^{[s],(a\omega)}_{m\lambda}(r) \notag\\
&\hspace{30mm}
 +\left(4\ii s\omega r-\lambda+s-a^2\omega^2+2am\omega\right)
        R^{[s],(a\omega)}_{m\lambda}(r)
        =0
\label{eq:radial-ode}
\end{align}
where $r\in (r_+,\infty)$.

The radial ODE \eqref{eq:radial-ode} has singularities at $r=r_{\pm}$, which are regular if $|a|<M$ but are a single irregular singularity of rank $1$ if $|a|=M$, and an irregular singularity of rank $1$ at $r=\infty$. By standard asymptotic analysis, we can span the solution space by two linearly independent asymptotic solutions at each of the singularities $r=r_+,\infty$. Our basis is given by the following:

\begin{definition}\label{def:radial-basis}
Fix $M>0$, $|a|\leq M$, $s\in \frac12\mathbb Z$ and an admissible frequency triple $(\omega,m,\lambda)$ with respect to $a$ and $s$.
\begin{enumerate}[label=\arabic*.]
\item Define $R^{[s]}_{\cH+}$ and $R^{[s]}_{\cH-}$ to be the unique classical solutions to the homogeneous radial ODE \eqref{eq:radial-ode} with boundary conditions
\begin{enumerate}[label=(\alph*)]
\item if $|a|<M$ and additionally $\omega-m\omega_+\neq 0$,
\begin{enumerate}[label=\roman*.]
\item $R^{[s]}_{\cH+}(r)(r-r_+)^{-\xi+s}$ and $R^{[s]}_{\cH-}(r)(r-r_+)^{\xi}$ are both smooth at $r=r_+$,
\item $\left((r-r_+)^{-\xi+s}R^{[s]}_{\cH+}\right)\big|_{r=r_+}=1$, and
 $\left((r-r_+)^{\xi}R^{[s]}_{\cH-}\right)\big|_{r=r_+}=1$;
\end{enumerate}
\item if $|a|<M$, $\omega=m\omega_+$, and $s\in \frac12\mathbb Z\backslash \mathbb Z_{> 0}$,
\begin{enumerate}[label=\roman*.]
\item $R^{[s]}_{\cH+}(r)(r-r_+)^{s}$ is smooth at $r=r_+$,
\item $\left((r-r_+)^{s}R^{[s]}_{\cH+}\right)\big|_{r=r_+}=1$,
\end{enumerate}
\item if $|a|<M$, $\omega=m\omega_+$, and $s\in \mathbb Z_{>0}$,
\begin{enumerate}[label=\roman*.]
\item $R^{[s]}_{\cH+}(r)$ is smooth at $r=r_+$,
\item $\left(R^{[s]}_{\cH+}\right)\big|_{r=r_+}=1$,
\end{enumerate}
\item if $|a|=M$ and additionally $\omega-m\omega_+\neq0$,
\begin{enumerate}[label=\roman*.]
\item $(r-M)^{2\ii M\omega+2s}R^{[s]}_{\cH+}(r)e^{-\beta(r-M)^{-1}}$ and $(r-M)^{-2\ii M\omega}R^{[s]}_{\cH-}(r)e^{\beta(r-M)^{-1}}$ are both smooth at $r=M$,
\item $\left((r-M)^{2\ii M\omega+2s}e^{-\beta(r-M)^{-1}}R^{[s]}_{\cH+}\right)\big|_{r=M}=1$, and
$\left((r-M)^{-2\ii M\omega}e^{\beta(r-M)^{-1}}R^{[s]}_{\cH-}\right)\big|_{r=M}=1$.
\end{enumerate}
\end{enumerate}
\item Define $R^{[s]}_{\cI+}$ and $R^{[s]}_{\cI-}$ to be the unique classical solutions to the homogeneous radial ODE \eqref{eq:radial-ode} and boundary conditions
\begin{enumerate}[label=(\alph*)]
\item $R^{[s]}_{\cI+}\sim e^{\ii\omega r}r^{2M\ii\omega-2s-1}$ and $R^{[s]}_{\cI-}\sim e^{-\ii\omega r}r^{-2M\ii\omega-1}$ asymptotically\footnote{We write $f\sim g$ asymptotically as $r\to\infty$ to mean that for every $N\geq 0$, we have complex numbers $c_1, \dots, c_N$ such that $f(r)=g(r)\left[1+\sum_{k=1}^N c_k r^{-k}+O(r^{-N-1})\right]$.} as $r\to\infty$,
\item $\left(e^{-\ii\omega r}r^{-2\ii M\omega+1+2s}R^{[s]}_{\cI+}\right)\big|_{r=\infty}=1$, and $\left(e^{\ii\omega r}r^{2\ii M\omega+1}R^{[s]}_{\cI-}\right)\big|_{r=\infty}=1$.
\end{enumerate}
\end{enumerate}
\end{definition}

In light of the previous comments, we find that we have the following representation for solutions of \eqref{eq:radial-ode}:

\begin{lemma}\label{lem:radial-representation}
Fix $M>0$, $|a|\leq M$, $s\in \frac12\mathbb Z$ and an admissible frequency triple $(\omega,m,\lambda)$ with respect to $a$ and $s$ such that $\omega\neq m\omega_+$. A solution $R^{[s],(a\omega)}_{m\lambda}$ to the homogeneous radial ODE \eqref{eq:radial-ode} can be written as, dropping most sub and superscripts,
\begin{equation}\label{eq:radial-representation}
        R^{[s]}=a^{[s]}_{\cH+}R^{[s]}_{\cH+}+a^{[s]}_{\cH-}R^{[s]}_{\cH-},
        \qquad
        R^{[s]}=a^{[s]}_{\cI+}R^{[s]}_{\cI+}+a^{[s]}_{\cI-}R^{[s]}_{\cI-},
\end{equation}
for some $a^{[s]}_{\cH\pm},a^{[s]}_{\cI\pm}\in\mathbb C$.
\end{lemma}

\subsubsection{The radial Teukolsky--Starobinsky identities}

In this section, we will introduce the Teukolsky--Starobinsky identities. We begin by defining the differential operators
\begin{equation}\label{eq:Dpm}
        D^{\pm}_n=\frac{\dd}{\dd r}\pm\ii\left(\frac{\omega(r^2+a^2)}{\Delta}-\frac{am}{\Delta}\right)+\frac{2n(r-M)}{\Delta}.
\end{equation}


The radial Teukolsky--Starobinsky identities are differential identities that relate solutions of the  radial ODE \eqref{eq:radial-ode} with spin $+s$ and spin $-s$, and were introduced in [TP74, SC74] and extended for general $s$ in [KMW89]. We will present a rigorous statement:

\begin{proposition}[Radial Teukolsky--Starobinsky identities]\label{prop:radial-ts}
Fix $M>0$, $|a|\leq M$ and $s\in \frac12\mathbb Z_{\geq0}$. Let $(\omega,m,\lambda)$ be an admissible frequency triple with respect to $s$ and $a$. Dropping most subscripts, let $R^{[\pm s]}$ be solutions to the homogeneous radial ODE \eqref{eq:radial-ode} of spin $\pm s$. If $\omega\neq m\omega_+$, consider the representation \eqref{eq:radial-representation}
\begin{align*}
R^{[\pm s]}&=a^{[\pm s]}_{\cH+}R^{[\pm s]}_{\cH+}
       +a^{[\pm s]}_{\cH-}R^{[\pm s]}_{\cH-} \\
&=a^{[\pm s]}_{\cI+}R^{[\pm s]}_{\cI+}
       +a^{[\pm s]}_{\cI-}R^{[\pm s]}_{\cI-};
\end{align*}
otherwise, we consider outgoing boundary conditions, i.e.\
\begin{align*}
R^{[\pm s]}&=a^{[\pm s]}_{\cH+}R^{[\pm s]}_{\cH+}\\
&=a^{[\pm s]}_{\cI+}R^{[\pm s]}_{\cI+},
\end{align*}
and set $a^{[\pm s]}_{\cH-}=0=a^{[\pm s]}_{\cI-}$ by convention. Then we have
\begin{align}
\Delta^s(D^+_0)^{2s}\left(\Delta^s R^{[+s]}\right)
&=\mathfrak C^{(1)}_s a^{[+s]}_{\cI+}R^{[-s]}_{\cI+}
  +\mathfrak C^{(7)}_s a^{[+s]}_{\cI-}R^{[-s]}_{\cI-} \notag\\
&=\mathfrak C^{(4)}_sa^{[+s]}_{\cH+}R^{[-s]}_{\cH+}
  +\mathfrak C^{(6)}_sa^{[+s]}_{\cH-}R^{[-s]}_{\cH-},
\notag\\[2mm]
\Delta^s(D^-_0)^{2s}R^{[-s]}
&=\mathfrak C^{(3)}_s a^{[-s]}_{\cI+}\Delta^sR^{[+s]}_{\cI+}
  +\mathfrak C^{(5)}_s a^{[-s]}_{\cI-}\Delta^sR^{[+s]}_{\cI-} \notag\\
&=\mathfrak C^{(2)}_sa^{[-s]}_{\cH+}\Delta^sR^{[+s]}_{\cH+}
  +\mathfrak C^{(8)}_sa^{[-s]}_{\cH-}\Delta^sR^{[+s]}_{\cH-},
\label{eq:TS-identities}
\end{align}
for some $\mathfrak C^{(i)}_s\in\mathbb C$, $i=1,\dots,8$, which depend only on $s$, $(a,M)$ and $(\omega,m,\lambda)$ and such that $\mathfrak C^{(i)}_0=1$ for $s=0$ and, for $s\neq 0$,
\begin{equation*}
        \mathfrak C^{(1)}_s=(2\ii\omega)^{2s},
        \qquad
        \mathfrak C^{(5)}_s=(-2\ii\omega)^{2s},
\end{equation*}
 and
\begin{equation}\label{eq:C2-C6}
\begin{gathered}
\mathfrak C^{(2)}_s=
\begin{cases}
\displaystyle \prod_{j=0}^{2s-1}(2\xi+s-j), & |a|<M, \text{~and either $\omega\neq m\omega_+$ or $s\in\frac12\mathbb Z_{\geq 0}\backslash\mathbb Z$},\\[2mm]
(-2\beta)^{2s}, & |a|=M,
\end{cases}
\\
\mathfrak C^{(6)}_s=
\begin{cases}
\displaystyle (r_+-r_-)^{2s}\prod_{j=0}^{2s-1}(-2\xi+s-j), & |a|<M, \text{~and either $\omega\neq m\omega_+$ or $s\in\frac12\mathbb Z_{\geq 0}\backslash\mathbb Z$},\\[2mm]
(2\beta)^{2s}, & |a|=M.
\end{cases}
\end{gathered}
\end{equation}
\end{proposition}

%\begin{remark}[Definition of the radial Teukolsky--Starobinsky constant]\label{rem:radial-ts-constant}
%Assume that $R^{[\pm s]}$ are outgoing solutions of the homogeneous radial ODE, i.e.\ $a^{[\pm s]}_{\cH-}=a^{[\pm s]}_{\cI-}=0$ in \eqref{eq:radial-representation}. Then, \eqref{eq:TS-identities} provide an alternative definition of the radial Teukolsky--Starobinsky constant from Proposition 2.10 as
%\begin{equation}\label{eq:Cs-def}
%        \mathfrak C_s:=\mathfrak C^{(1)}_s\mathfrak C^{(3)}_s
%        =\mathfrak C^{(2)}_s\mathfrak C^{(4)}_s.
%\end{equation}
%If $R^{[\pm s]}$ are not outgoing, we consider the radial Teukolsky--Starobinsky constant as defined in Proposition 2.10 and conclude
%\begin{equation*}
%        \mathfrak C^{(1)}_s\mathfrak C^{(3)}_s
%        =\mathfrak C^{(2)}_s\mathfrak C^{(4)}_s
%        =\mathfrak C^{(5)}_s\mathfrak C^{(7)}_s
%        =\mathfrak C^{(6)}_s\mathfrak C^{(8)}_s
%        =\mathfrak C_s.
%\end{equation*}
%Having defined the radial Teukolsky--Starobinsky constant, we can write \eqref{eq:TS-identities} in terms of this constant and the $\mathfrak C^{(i)}_s$ for which we have simple formulas in terms of the frequency parameters:
%\begin{align}
%\Delta^s(D^+_0)^{2s}\left(\Delta^sR^{[+s]}\right)
%&=\mathfrak C^{(1)}_s a^{[+s]}_{\cI+}R^{[-s]}_{\cI+}
%  +\mathfrak C_s\left(\mathfrak C^{(5)}_s\right)^{-1}a^{[+s]}_{\cI-}R^{[-s]}_{\cI-} \notag\\
%&=\mathfrak C_s\left(\mathfrak C^{(2)}_s\right)^{-1}a^{[+s]}_{\cH+}R^{[-s]}_{\cH+}
%  +\mathfrak C^{(6)}_s a^{[+s]}_{\cH-}R^{[-s]}_{\cH-},
%\notag\\[2mm]
%\Delta^s(D^-_0)^{2s}R^{[-s]}
%&=\mathfrak C_s\left(\mathfrak C^{(1)}_s\right)^{-1}a^{[-s]}_{\cI+}\Delta^sR^{[+s]}_{\cI+}
%  +\mathfrak C^{(5)}_s a^{[-s]}_{\cI-}\Delta^sR^{[+s]}_{\cI-} \notag\\
%&=\mathfrak C^{(2)}_s a^{[-s]}_{\cH+}\Delta^sR^{[+s]}_{\cH+}
%  +\mathfrak C_s\left(\mathfrak C^{(6)}_s\right)^{-1}a^{[-s]}_{\cH-}\Delta^sR^{[+s]}_{\cH-}.
%\label{eq:TS-identities-C}
%\end{align}
%\end{remark}

In order to prove Proposition \ref{prop:radial-ts}, we recall the following two results from [KMW89]:

\begin{lemma}\label{lem:P-ode} For solutions of the homogeneous radial ODE~\eqref{eq:radial-ode} with spin $\pm s$, we set
\begin{align*}
P^{[+ s]}:=\Delta^s R^{[+s]}, \quad P^{[-s]}:=R^{[-s]}.
\end{align*}
Then, we have
\begin{equation}\label{eq:P-ode}
\left[\Delta D^{\mp}_{1-s}D^{\pm}_0\pm2(2s-1)\ii\omega r\right]P^{[\pm s]}(r)
 =\left(\lambda+|s|+a^2\omega^2-2am\omega\right)P^{[\pm s]}(r).
\end{equation}
\end{lemma}

By induction on $2s$, one can also show that:

\begin{lemma}\label{lem:radial-commutation}
The operators $D^{\pm}_0$, defined in \eqref{eq:Dpm}, satisfy the commutation relation
\begin{align}\label{eq:radial-commutation}
&\Delta^s(D^{\mp}_0)^{2s}
\left[\Delta D^{\pm}_{1-s}D^{\mp}_0\mp2(2s-1)\ii\omega r\right]
\notag\\
&\hspace{35mm}=
\left[\Delta D^{\mp}_{1-s}D^{\pm}_0\pm2(2s-1)\ii\omega r\right]
\Delta^s(D^{\mp}_0)^{2s}.
\end{align}
\end{lemma}

Combining the two lemmas, we see that, if $P^{[\pm s]}$ solves \eqref{eq:P-ode} with spin $\pm$, then $\Delta^s(D^{\pm}_0)^{2s}P^{[\pm s]}$ is a solution of the homogeneous radial ODE \eqref{eq:P-ode} for spin $\mp s$. We are now ready to prove Proposition \ref{prop:radial-ts}:

\begin{proof}[Proof of Proposition \ref{prop:radial-ts}]
The statement is trivial for $s=0$. For $s\neq0$, we obtain the result by using the asymptotic representations of $R^{[\pm s]}_{\cH+}$, $R^{[\pm s]}_{\cI+}$, $R^{[\pm s]}_{\cH-}$ and $R^{[\pm s]}_{\cI-}$, on which we can act directly with the Teukolsky--Starobinsky operators. Since $\omega\neq 0$, these are
\begin{align}
\Delta^{\frac{s\pm s}{2}}R^{[\pm s]}_{\cI+}
&=e^{\ii\omega r}r^{2\ii M\omega+s\mp s-1}
\left[\sum_{k=0}^{2s}c^{[\pm s]}_kr^{-k}+O(r^{-2s-1})\right],
\notag\\
\Delta^{\frac{s\pm s}{2}}R^{[\pm s]}_{\cI-}
&=e^{-\ii\omega r}r^{-2\ii M\omega+s\pm s-1}
\left[\sum_{k=0}^{2s}c^{[\pm s]}_{k,2}r^{-k}+O(r^{-2s-1})\right],
\label{eq:asymptotic-representations-infty}
\end{align}
as $r\to \infty$; if $\omega\neq m\omega_+$,
\begin{align}
\Delta^{\frac{s\pm s}{2}}R^{[\pm s]}_{\cH+}
&=
\begin{cases}
\displaystyle (r-r_-)^{\frac{s\pm s}{2}}\sum_{k=0}^{\infty}b^{[\pm s]}_k(r-r_+)^{\xi+(s\mp s)/2+k}, & |a|<M,\\[2mm]
\displaystyle \sum_{k=0}^{\infty}b^{[\pm s]}_ke^{\beta(r-M)^{-1}}(r-M)^{-2\ii M\omega+s\mp s+k}, & |a|=M,
\end{cases}
\notag\\
\Delta^{\frac{s\pm s}{2}}R^{[\pm s]}_{\cH-}
&=
\begin{cases}
\displaystyle (r-r_-)^{\frac{s\pm s}{2}}\sum_{k=0}^{\infty}b^{[\pm s]}_{k,2}(r-r_+)^{-\xi+(s\pm s)/2+k}, & |a|<M,\\[2mm]
\displaystyle \sum_{k=0}^{\infty}b^{[\pm s]}_{k,2}e^{-\beta(r-M)^{-1}}(r-M)^{2\ii M\omega+s\pm s+k}, & |a|=M,
\end{cases}
\label{eq:asymptotic-representations-hor1}
\end{align}
as $r\to r_+$; 
and, if $\omega=m\omega_+$,
\begin{align}
\Delta^{\frac{s\pm s}{2}}R^{[\pm s]}_{\cH+}
&= (r-r_-)^{\frac{s\pm s}{2}}\sum_{k=0}^{\infty}b^{[\pm s]}_k(r-r_+)^{\xi+ (s\mp s)/2+k}, & |a|<M, \, s\in\frac12\mathbb Z_{\geq 0}\backslash \mathbb Z \\
R^{[- s]}_{\cH+}
&= \sum_{k=0}^{\infty}b^{[- s]}_k(r-r_+)^{\xi+ s+k}, & |a|<M,\, s\in \mathbb Z_{\geq 0},\\
\Delta^s R^{[+s]}_{\cH+}
&=
\displaystyle (r-r_-)^s\sum_{k=0}^{\infty}b^{[+s]}_k(r-r_+)^{\xi+s+k}, & |a|<M, \, s\in \mathbb Z_{\geq 0},
\label{eq:asymptotic-representations-hor2}
\end{align}
as $r\to r_+$.
 By the normalizations, $c^{[\pm s]}_0=c^{[\pm s]}_{0,2}=b^{[\pm s]}_0=b^{[\pm s]}_{0,2}=1$.

For computations at $r=r_+$, it will be useful to recall from the definition of $\xi$ and $\beta$ in \eqref{eq:xi-beta},
\begin{align}
\xi(r-r_-)+\ii\omega(r^2+a^2)-\ii am
&=\xi(r_+-r_-)+\ii\omega(r_+^2+a^2)-\ii am
\notag\\
&\quad+\xi(r-r_+)+\ii\omega(r-r_+)(r+r_+)
\notag\\
&=\xi(r-r_+)+\ii\omega(r-r_+)(r+r_+),
\label{eq:xi-identity}\\
\beta+\ii\omega(r-M)(r+M)&=\ii\omega(r^2+M^2)-\ii am.
\label{eq:beta-identity}
\end{align}

\emph{The cases $R^{[+s]}_{\cI+}$ and $R^{[-s]}_{\cH+}$.} In these cases, the constant arising from application of the Teukolsky--Starobinsky operators to these functions involves only the first term of the series.

We begin at $r\to\infty$, with positive spin. For any $s\geq\frac12$, we will have
\begin{align*}
D^+_0\left(\Delta^s R^{[+s]}_{\cI+}\right)
&=\left(\frac{\dd}{\dd r}+\ii\frac{\omega(r^2+a^2)-am}{\Delta}\right)
        \left(\Delta^sR^{[+s]}_{\cI+}\right)\\
&=e^{\ii\omega r}r^{2\ii M\omega-1}
\left[\sum_{k=0}^{2s}c^{[+s]}_k
\left(2\ii\omega-\frac{k+1-2\ii M\omega}{r}
      +\frac{2M\ii\omega r-\ii am}{\Delta}\right)r^{-k}
      +O(r^{-2s-1})\right].
\end{align*}
Repeating $2s-1$ times, we obtain
\begin{equation}\label{eq:Iplus-easy}
\Delta^s(D^+_0)^{2s}\left(\Delta^sR^{[+s]}_{\cI+}\right)
=e^{\ii\omega r}r^{2\ii M\omega+2s-1}
\left[\sum_{k=0}^{2s}c^{[+s]}_k\left((2\ii\omega)^{2s}+O(r^{-1})\right)r^{-k}+O(r^{-2s-1})\right].
\end{equation}
By Lemma \ref{lem:radial-commutation}, \eqref{eq:Iplus-easy} satisfies the spin $-s$ radial ODE, hence, by Definition \ref{def:radial-basis}, we conclude
\begin{equation*}
\Delta^s(D^+_0)^{2s}\left(\Delta^sR^{[+s]}_{\cI+}\right)=\mathfrak C^{(1)}_sR^{[-s]}_{\cI+},
\qquad
\mathfrak C^{(1)}_s:=(2\ii\omega)^{2s}.
\end{equation*}

We now turn to $r=r_+$, and consider solutions with spin $-s$. We start with $|a|<M$; for any $s\neq 0$, using \eqref{eq:xi-identity}, we obtain
\begin{align}
D^-_0R^{[-s]}_{\cH+}
&=\left(\frac{\dd}{\dd r}-\ii\frac{\omega(r^2+a^2)}{\Delta}+\frac{\ii am}{\Delta}\right)R^{[-s]}\notag\\
&=\sum_{k=0}^{\infty}b^{[-s]}_k
\left((\xi+k+s)-\frac{\ii\omega(r^2+a^2)}{r-r_-}+\frac{\ii am}{r-r_-}\right)(r-r_+)^{\xi+k+s-1}\notag\\
&=\sum_{k=0}^{\infty}b^{[-s]}_k
\left((2\xi+k+s)-\frac{(\xi+\ii\omega(r+r_+))(r-r_+)}{r-r_-}\right)(r-r_+)^{\xi+k+s-1}. \label{eq:Hplus-easy-subextremal-int}
\end{align}
If $s\in\mathbb Z_{\geq 0}$, repeating the procedure $s-1$ more times yields 
\begin{equation*}
(D^-_0)^{s}R^{[-s]}_{\cH+}
=\sum_{k=0}^{\infty}b^{[-s]}_k
\left[\prod_{j=0}^{s-1}(2\xi+k+s-j)+O(r-r_+)\right]
(r-r_+)^{\xi+k},
\end{equation*}
which is clearly regular at $r=r_+$ if $\omega=m\omega_+$; this regularity is preserved by acting with a further factor of $(D^-_0)^{s}$: 
\begin{equation*}
\Delta^s(D^-_0)^{2s}R^{[-s]}_{\cH+}
=\Delta^s\left[s! f_-(b_0^{[-s]},\dots, b_s^{[-s]}, \omega=m\omega_+, a , M)+O(r-r_+)\right],\label{eq:Hplus-easy-subextremal1}
\end{equation*}
for some function $f_-$ which is independent of $r$. 
If, on the other hand, $s\in\frac12\mathbb Z_{\geq 0}\backslash\mathbb Z$ or if $\omega\neq m\omega_+$, repeating the procedure of \eqref{eq:Hplus-easy-subextremal-int} $2s-1$ more times yields
\begin{equation}\label{eq:Hplus-easy-subextremal2}
\Delta^s(D^-_0)^{2s}R^{[-s]}_{\cH+}
=\sum_{k=0}^{\infty}b^{[-s]}_k
\left[\prod_{j=0}^{2s-1}(2\xi+k+s-j)+O(r-r_+)\right]
\Delta^s(r-r_+)^{\xi+k-s}.
\end{equation}
Finally, if $|a|=M$, then for any $s\neq 0$ and $\omega\neq m\omega_+$, using \eqref{eq:beta-identity}, we obtain
\begin{align*}
D^-_0R^{[-s]}_{\cH+}
&=\left(\frac{\dd}{\dd r}-\ii\omega\frac{r^2+M^2}{(r-M)^2}+\frac{\ii am}{(r-M)^2}\right)R^{[-s]}\\
&=\left(\frac{\dd}{\dd r}-\frac{\beta}{(r-M)^2}-\frac{(r+M)\ii\omega}{r-M}\right)
   e^{\beta(r-M)^{-1}}\sum_{k=0}^{\infty}b^{[-s]}_k(r-M)^{-2\ii M\omega+2s+k}\\
&=\left[-2\beta+O(r-M)\right]e^{\beta(r-M)^{-1}}
   \sum_{k=0}^{\infty}b^{[-s]}_k(r-M)^{-2\ii M\omega+2s-2+k},
\end{align*}
which, repeating $2s-1$ times, yields
\begin{equation}\label{eq:Hplus-easy-extremal}
\Delta^s(D^-_0)^{2s}R^{[-s]}_{\cH+}
=\sum_{k=0}^{\infty}b^{[-s]}_k\left[(-2\beta)^{2s}+O(r-M)\right]
 e^{\beta(r-M)^{-1}}(r-M)^{-2\ii M\omega+k}.
\end{equation}
By Lemma \ref{lem:radial-commutation}, $\Delta^s(D^-_0)^{2s}R^{[-s]}_{\cH+}$ satisfies the homogeneous spin $+s$ radial ODE \eqref{eq:P-ode}, hence, by Definition \ref{def:radial-basis}, we conclude that
\begin{equation*}
\Delta^s(D^-_0)^{2s}R^{[-s]}_{\cH+}=\mathfrak C^{(2)}_s\Delta^sR^{[+s]}_{\cH+},
\end{equation*}
where we have set
\begin{gather*}
\mathfrak C^{(2)}_s:=
\begin{dcases}
\prod_{j=0}^{2s-1}(2\xi+s-j)\quad&\text{if $|a|<M$ and either $\omega\neq m\omega_+$ or $s\in\frac12\mathbb Z_{\geq 0}\backslash\mathbb Z$}, \\
s! f_-(b_0^{[-s]},\dots, b_s^{[-s]}, \omega, a , M),&\text{if $|a|<M$, $\omega=m\omega_+$ and $s\in\mathbb Z_{\geq 0}$},\\ 
(-2\beta)^{2s}\quad&\text{if }|a|=M.
\end{dcases}
\end{gather*}

\emph{The cases $R^{[+s]}_{\cH+}$ and $R^{[-s]}_{\cI+}$.} We need to take the same series-based approach as in the previous case. As we will see, for $R^{[+s]}_{\cH+}$ and $R^{[-s]}_{\cI+}$, the constant arising from application of the Teukolsky--Starobinsky operators to the series for these functions is more complicated to compute, as it involves up to $2s$ terms of their series expansion.

We begin with spin $-s$ at $r=\infty$. For any $s\geq\frac12$, we have
\begin{align}
D^-_0R^{[-s]}_{\cI+}
&=\left(\frac{\dd}{\dd r}-\ii\frac{\omega(r^2+a^2)}{\Delta}+\frac{\ii am}{\Delta}\right)
\left[e^{\ii\omega r}r^{2\ii M\omega+2s-1}
\sum_{k=0}^{2s}c^{[-s]}_kr^{-k}+O(r^{-2s-1})\right]
\notag\\
&=e^{\ii\omega r}r^{2\ii M\omega+2s-1}
\left[
\sum_{k=0}^{2s}c^{[-s]}_k\left(\frac{2\ii M\omega (a^2-2M r)}{r\Delta}+\frac{\ii am}{\Delta}+\frac{2s-1-k}{r}\right)r^{-k}\right].
\label{eq:Dminus-Iplus}
\end{align}
The key difficulty, which is not present in the case of spin $+s$, is the cancellation that occurs between the $\omega$ dependence of $D^-_0$ and that of the leading order behavior of $R^{[-s]}_{\cI+}$. For instance, if $s=\frac12$, then
\begin{equation*}
D^-_0R^{[-1/2]}_{\cI+}
=\left[(-4\ii M^2\omega+\ii am)c^{[-1/2]}_0-c^{[-1/2]}_1+O(r^{-1})\right]
 e^{\ii\omega r}r^{2\ii M\omega-2},
\end{equation*}
so taking $D^-_0$ lowers the leading order behavior by a factor of $r^{-2}$. Indeed, an application of $(D^-_0)^{2s}$ lowers the leading order polynomial decay of $R^{[-s]}_{\cI+}$ by a factor of $r^{-4s}$ and the coefficient of the leading order term on the right hand side depends on coefficients $c^{[-s]}_0,\dots,c^{[-s]}_{2s}$ of the asymptotic expansion for $R^{[-s]}_{\cI+}$:
\begin{equation}\label{eq:C3-asymptotic}
\Delta^s(D^-_0)^{2s}R^{[-s]}_{\cI+}
=\Delta^s e^{\ii\omega r}r^{2\ii M\omega-2s-1}
\left[f_+\left(c^{[-s]}_0,\dots,c^{[-s]}_{2s},\omega,m,a, M\right)+O(r^{-1})\right],
\end{equation}
where $f_+$ is independent of $r$. By standard theory of asymptotic analysis, $c^{[-s]}_1,\dots,c^{[-s]}_{2s}$ are uniquely determined by $c^{[-s]}_0$ and the ODE parameters $(\omega,m,\lambda,s,a, M)$, so we can define
\begin{equation*}
\mathfrak C^{(3)}_s(\omega,m,\lambda,s,a, M):=
\frac{f_+\left(c^{[-s]}_0,\dots,c^{[-s]}_{2s},\omega,m,a, M\right)}{c^{[-s]}_0}.
\end{equation*}
By Lemma \ref{lem:radial-commutation}, \eqref{eq:C3-asymptotic} satisfies the spin $+s$ radial ODE. As before, by Definition \ref{def:radial-basis}, we conclude
\begin{equation*}
\Delta^s(D^-_0)^{2s}R^{[-s]}_{\cI+}=\mathfrak C^{(3)}_s\Delta^sR^{[+s]}_{\cI+}.
\end{equation*}

Now consider spin $+s$ near $r=r_+$, and recall \eqref{eq:xi-identity} and \eqref{eq:beta-identity}. If $|a|<M$, and $\omega\neq m\omega_+$ or $s\in\frac12\mathbb Z_{\geq 0}\backslash\mathbb Z$, then
\begin{align*}
D^+_0\left(\Delta^sR^{[+s]}_{\cH+}\right)
&=\left(\frac{\dd}{\dd r}+\frac{\ii\omega(r+r_+)}{r-r_-}-\frac{\xi}{r-r_+}+\frac{\xi}{r-r_-}\right)
\sum_{k=0}^{\infty}b^{[+s]}_k(r-r_+)^{\xi+k}(r-r_-)^s\\
&=\sum_{k=0}^{\infty}b^{[+s]}_k
\left(\frac{k}{r-r_+}+\frac{\xi+s+\ii\omega(r+r_+)}{r-r_-}\right)
(r-r_+)^{\xi+k}(r-r_-)^s,
\end{align*}
and if $|a|=M$,
\begin{align*}
D^+_0\left(\Delta^sR^{[+s]}_{\cH+}\right)
&=\left(\frac{\dd}{\dd r}+\frac{\beta}{(r-M)^2}+\frac{2M\ii\omega}{r-M}+\ii\omega\right)
\sum_{k=0}^{\infty}b^{[+s]}_ke^{\beta(r-M)^{-1}}(r-M)^{-2\ii M\omega+k}\\
&=\sum_{k=0}^{\infty}b^{[+s]}_k\left(\frac{k}{r-M}+\ii\omega\right)
 e^{\beta(r-M)^{-1}}(r-M)^{-2\ii M\omega+k}.
\end{align*}
As in the previous case, there is cancellation between the $\xi$ or $\beta$ dependence of the Teukolsky--Starobinsky operators and that of the leading order behavior of $\Delta^sR^{[+s]}_{\cH+}$. Proceeding as above, we note that application of $(D^+_0)^{2s}$ does not change the leading order polynomial decay of $\Delta^s R^{[+s]}_{\cH+}$ as $r\to r_+$ if $\omega\neq m\omega_+$ or, in case $|a|<M$, if $s\in\frac12\mathbb Z\backslash \mathbb Z$, but makes the coefficient of the resulting leading order term depend on coefficients $b^{[+s]}_0,\dots,b^{[+s]}_{2s}$ of the asymptotic series for $R^{[+s]}_{\cH+}$. In the case $|a|<M$, $s\in\mathbb Z$ and $\omega=m\omega_+$, we have
\begin{align*}
D^+_0\left(\Delta^sR^{[+s]}_{\cH+}\right)
&=\left(\frac{\dd}{\dd r}+\frac{\ii\omega(r+r_+)}{r-r_-}\right)
\sum_{k=0}^{\infty}b^{[+s]}_k(r-r_+)^{k}(r-r_-)^s\\
&=\sum_{k=0}^{\infty}b^{[+s]}_k\left[s+k+O(r-r_+)\right](r-r_+)^{\xi+s+k-1}(r-r_-)^s\\
\implies (D^+_0)^s\left(\Delta^sR^{[+s]}_{\cH+}\right)
&=\sum_{k=0}^{\infty}b^{[+s]}_k\left[\prod_{j=0}^{s-1}(s+k-j)+O(r-r_+)\right](r-r_+)^{k}(r-r_-)^s\\
\implies \Delta^s (D^+_0)^{2s}\left(\Delta^sR^{[+s]}_{\cH+}\right)
&=\left[s! g_-(b_0^{[-s]},\dots, b_s^{[-s]}, \omega=m\omega_+, a , M)+O(r-r_+)\right](r-r_+)^s,
\end{align*}
for some function $g_-$ which is independent of $r$. Thus, we again appeal to theory of asymptotic analysis to define $\mathfrak C^{(4)}_s(\omega,m,\lambda,s,a,M)$ such that
\begin{equation*}
\Delta^s(D^+_0)^{2s}\left(\Delta^sR^{[+s]}_{\cH+}\right)
=\mathfrak C^{(4)}_sR^{[-s]}_{\cH+}.
\end{equation*}

\emph{The cases $R^{[\pm s]}_{\cH-}$ and $R^{[\pm s]}_{\cI-}$.} Finally by direct inspection of \eqref{eq:asymptotic-representations-infty}--\eqref{eq:asymptotic-representations-hor2}, for $R^{[\pm s]}_{\cH-}$ and $R^{[\pm s]}_{\cI-}$, we can easily obtain the Teukolsky--Starobinsky identities by analogy with $R^{[\mp s]}_{\cH+}$ and $R^{[\mp s]}_{\cI+}$, respectively.
\end{proof}


\subsection{The angular case}

\subsubsection{The angular ODE and its solutions}

Fix parameters $s\in \frac12\Z_{\geq0}$, $m-s\in\Z$, $\nu\in\R$ and $\lambda\in\R$.

In this section, we discuss the angular ODE
\begin{equation}\label{eq:angular-ode}
-\frac{1}{\sin\theta}\frac{d}{d\theta}\left(\sin\theta\frac{d}{d\theta}\right)\Xi^{[s]}
+\left(\frac{(m+ s\cos\theta)^2}{\sin^2\theta}
+\nu^2\sin^2\theta+2\nu s\cos\theta\right)\Xi^{[s]}
=(\lambda+\nu^2)\Xi^{[s]},
\end{equation}
where $\theta\in(0,\pi)$. It will be convenient to sometimes use the variable $x=\cos\theta$, so that \eqref{eq:angular-ode} is considered for $x\in(-1,1)$.

The ODE has regular singularities at $\theta=0,\pi$ ($x=\pm 1$), and an irregular singularity at $\theta=\infty$. By standard asymptotic analysis, we can span the solution space by two linearly independent asymptotic solutions at each of the singularities $\theta=0,\pi$. Our basis is given by the following:

\begin{definition}\label{def:angular-basis}
Fix $s,m\in \frac12\mathbb Z$ with $m-s\in \mathbb Z$, and $(\lambda,\nu)\in\mathbb R^2$.
\begin{enumerate}[label=\arabic*.]
\item Define $\Xi^{[s]}_{\cN^+}$ to be the unique classical solution to the angular ODE \eqref{eq:angular-ode} such that 
\begin{enumerate}[label=\roman*.]
\item $\Xi^{[s]}_{\cN^+}(\theta)(1-\cos\theta)^{-\frac12|m+s|}$ is smooth at $\theta=0$,
\item $\left((1-\cos\theta)^{-\frac12|m+s|}\Xi^{[s]}_{\cN^+}\right)\big|_{\theta=0}=1$.
\end{enumerate}
\item Define $\Xi^{[s]}_{\cS^+}$ to be the unique classical solution to the angular ODE \eqref{eq:angular-ode} such that 
\begin{enumerate}[label=\roman*.]
\item $\Xi^{[s]}_{\cS^+}(\theta)(1+\cos\theta)^{-\frac12|m-s|}$ is smooth at $\theta=\pi$,
\item $\left((1+\cos\theta)^{-\frac12|m-s|}\Xi^{[s]}_{\cS^+}\right)\big|_{\theta=\pi}=1$.
\end{enumerate}
\end{enumerate}
\end{definition}

In light of the previous comments, we find that we have the following representation for solutions of \eqref{eq:angular-ode}:

\begin{lemma}\label{lem:angular-representation}
 [TO FILL]
\end{lemma}

\subsubsection{The angular Teukolsky--Starobinsky identities}

In this section, we will introduce the Teukolsky--Starobinsky identities.  We begin by defining the differential operators
\begin{equation}\label{eq:Lpm}
        L^{\pm}_n=\frac{\dd}{\dd\theta}\pm\left(\frac{m}{\sin\theta}-\nu\sin\theta\right)+n\cot\theta.
\end{equation}

The angular Teukolsky--Starobinsky identities are differential identities that relate solutions of the  angular ODE \eqref{eq:angular-ode} with spin $+s$ and spin $-s$, and were introduced in [TP74, SC74] and extended for general $s$ in [KMW89]. We will present a rigorous statement:

\begin{proposition}[Angular Teukolsky--Starobinsky identities]\label{prop:angular-ts} Fix $s\in \frac12\mathbb Z_{\geq0}$, $m-s\in\mathbb Z$ and $(\lambda,\nu)\in \mathbb R^2$. Consider regular solutions to the angular ODE \eqref{eq:angular-ode}, i.e.\ solutions that admit a representation
\begin{align*}
\Xi^{[\pm s]}&=a^{[\pm s]}_{\cN^+}\Xi^{[\pm s]}_{\cN^+}\\
&=a^{[\pm s]}_{\cS^+}\Xi^{[\pm s]}_{\cS^+},
\end{align*}
for some complex constants $a_{\cN^+}^{[\pm s]}$, $a_{\cS^+}^{[\pm s]}$. Then we have
\begin{align}
(\sin\theta)^{2s}\left(\frac{L^{+}_s}{\sin\theta}\right)^{2s}\Xi^{[+s]}
&=\mathfrak B^{(1)}_s a^{[+s]}_{\cN^+}\Xi^{[-s]}_{\cN^+} \notag\\
&=\mathfrak B^{(4)}_sa^{[+s]}_{\cS^+}\Xi^{[-s]}_{\cS^+},
\notag\\[2mm]
(\sin\theta)^{2s}\left(\frac{L^{-}_s}{\sin\theta}\right)^{2s}\Xi^{[-s]}
&=\mathfrak B^{(3)}_s a^{[-s]}_{\cN^+}\Xi^{[+s]}_{\cN^+}
\notag\\
&=\mathfrak B^{(2)}_sa^{[-s]}_{\cS^+}\Xi^{[+s]}_{\cS^+},
\label{eq:angular-TS-identities}
\end{align}
for some $\mathfrak B^{(i)}_s\in\mathbb C\backslash\{0\}$, $i=1,\dots,4$, which depend only on $s$, $m$ and $(\lambda,\nu)$ and such that $\mathfrak B^{(i)}_0=1$ for $s=0$ and, for $s\neq 0$,
\begin{align*}
 \mathfrak B^{(1)}_s&=\dots,\\
 \mathfrak B^{(2)}_s&=(-1)^{2s} \mathfrak B^{(1)}_s,\\
 \mathfrak B^{(3)}_s&=\dots,\\
 \mathfrak B^{(4)}_s&=(-1)^{2s} \mathfrak B^{(3)}_s.
\end{align*}
\end{proposition}

In order to prove Proposition \ref{prop:angular-ts}, we recall the following two results from [KMW89]:

\begin{lemma}\label{lem:angular-ode-L} For solutions of the angular ODE~\eqref{eq:angular-ode} with spin $\pm s$,  we have
\begin{equation}\label{eq:angular-ode-L}
\left[L^{\mp}_{1-s}L^{\pm}_s\mp2(2s-1)\nu\cos\theta\right]\Xi^{[\pm s]}=-(\lambda+\nu^2-2m\nu+|s|)\Xi^{[\pm s]}
\end{equation}
\end{lemma}

By induction on $2s$, one can also show that:

\begin{lemma}\label{lem:angular-commutation}
The operators $L^{\pm}_s$, defined in \eqref{eq:Lpm}, satisfy the commutation relation
%\begin{align}\label{eq:angular-commutation}
%&(\sin\theta)^s\left(\frac{L^{\mp}_s}{\sin\theta}\right)^{2s}
%\left[ L^{\pm}_{1-s}L^{\mp}_s\pm2(2s-1)\nu\cos\theta\right]
%\notag\\
%&\hspace{35mm}=
%\left[L^{\mp}_{1-s}L^{\pm}_s\mp2(2s-1)\nu\cos\theta\right]
%(\sin\theta)^s\left(\frac{L^{\mp}_s}{\sin\theta}\right)^{2s}.
%\end{align}
\begin{align}\label{eq:angular-commutation}
&(\sin\theta)^{2s}\left(\frac{L^{\pm}_s}{\sin\theta}\right)^{2s}
\left[ L^{\mp}_{1-s}L^{\pm}_s\mp2(2s-1)\nu\cos\theta\right]
\notag\\
&\hspace{35mm}=
\left[L^{\pm}_{1-s}L^{\mp}_s\pm2(2s-1)\nu\cos\theta\right]
(\sin\theta)^{2s}\left(\frac{L^{\pm}_s}{\sin\theta}\right)^{2s}.
\end{align}
\end{lemma}

Combining the two lemmas, we see that, if  $\Xi^{[\pm s]}$ solves \eqref{eq:angular-ode-L} for spin $\pm s$, then $(\sin\theta)^{2s}\left(\frac{L^{\pm}_s}{\sin\theta}\right)^{2s}\Xi^{[\pm s]}$ is a solution to the angular ODE \eqref{eq:angular-ode-L} with spin $\mp s$. We are now ready to prove Proposition \ref{prop:angular-ts}:

\begin{proof}[Proof of Proposition \ref{prop:angular-ts}]
The statement is trivial for $s=0$. For $s\neq 0$, we obtain the result by using the asymptotic representations of $\Xi_{\cN^+}$ and $\Xi_{\cS^+}$: with respect to the variable $x=\cos\theta$,
\begin{align*}
\Xi^{[\pm s]}_{\cN^+}&=\sum_{k=0}^\infty b_k^{[\pm s]}(1-x)^{\frac12|m\pm s|+k},\qquad 
\Xi^{[\pm s]}_{\cS^+}=\sum_{k=0}^\infty c_k^{[\pm s]}(1+x)^{\frac12|m\mp s|+k},
\end{align*}
on which we can act directly with the Teukolsky--Starobinsky operators. By the normalizations, $c_0^{[\pm s]}=b_0^{[\pm s]}=1$. Note also that 
\begin{align}
\frac{L^{\pm}_s}{\sin\theta}=-\frac{d}{dx}+\frac{sx\pm m}{1-x^2}\mp \nu= -\frac{d}{dx}+\frac{s\pm m}{2(1-x)}-\frac{s\mp m}{2(1+x)}\mp \nu. \label{eq:angular-TS-op-rewritten}
\end{align}

\medskip
\noindent \textit{The $\theta=0$ end.}
We begin by computing the effect of applying the Teukolsky--Starobinsky operators to $\Xi^{[\pm s]}_{\cN^+}$. In doing so, it will be useful to note that, from \eqref{eq:angular-TS-op-rewritten},
\begin{align*}
\left(\frac{L^{\pm}_s}{\sin\theta}\right)\left((1-x)^{-\frac{1}{2}(s\pm m)}\Xi\right) = (1-x)^{-\frac{1}{2}(s\pm m)}\left(-\frac{d}{dx}-\frac{s\mp m}{2(1+x)}\mp \nu\right)\Xi,
\end{align*}
for any function $\Xi(x)$. It follows immediately that, if $s\pm m\leq 0$, then
\begin{align*}
(\sin\theta)^{2s}\left(\frac{L^{\pm}_s}{\sin\theta}\right)^{2s} \Xi^{[\pm s]}_{\cN^+} 
&=(\sin\theta)^{2s}\left(\frac{L^{\pm}_s}{\sin\theta}\right)^{2s} \left[(1-x)^{-\frac12(s\pm m)}\sum_{k=0}^\infty b_k^{[\pm s]}(1-x)^{k}\right]\\
&=
(1+x)^{s}(1-x)^{\frac12|-s\pm m|}\left(-\frac{d}{dx}-\frac{s\mp m}{2(1+x)}\mp \nu\right)^{2s} \sum_{k=0}^\infty b_k^{[\pm s]}(1-x)^{k}.
\end{align*} 
Combining Lemmas~\ref{lem:angular-commutation}, \ref{lem:angular-ode-L} and Definition~\ref{def:angular-basis}, we deduce that $(\sin\theta)^{2s}\left(\frac{L^{\pm}_s}{\sin\theta}\right)^{2s} \Xi^{[\pm s]}_{\cN^+} $ is a solution to the $\mp s$ angular ODE with no irregular behavior at $\theta= 0$. Thus, there are functions $f_\pm$ so that
\begin{gather*}
s+ m \leq 0\implies \mathfrak B_s^{(1)}= 2^{s}\left((2s)!b_{2s}^{[+s]}+f_+(b_0^{[+s]}, \dots, b_{2s-1}^{[+s]})\right), \\ s- m \leq 0\implies \mathfrak B_s^{(3)}=2^{s}\left((2s)!b_{2s}^{[-s]}+f_-(b_0^{[-s]}, \dots, b_{2s-1}^{[-s]})\right).
\end{gather*}

If $s\pm m\geq 1$, then
\begin{align*}
\left(\frac{L^{\pm}_s}{\sin\theta}\right) \Xi^{[\pm s]}_{\cN^+} 
&=\sum_{k=0}^\infty b_k^{[\pm s]}\left[k+s\pm m+O(1-x)\right](1-x)^{\frac12(s\pm m-2)+k}\\
\implies \left(\frac{L^{\pm}_s}{\sin\theta}\right)^n \Xi^{[\pm s]}_{\cN^+} &=\sum_{k=0}^\infty b_k^{[\pm s]}\left[\prod_{j=0}^{n-1} (k+s\pm m-j)+O(1-x)\right](1-x)^{\frac12(s\pm m-2n)+k}.
\end{align*}
It follows that:
\begin{itemize}
\item if $s\pm m \geq 2s\geq 1$, then 
\begin{align*}
&(\sin\theta)^{2s}\left(\frac{L^{\pm}_s}{\sin\theta}\right)^{2s} \Xi^{[\pm s]}_{\cN^+}  \\
%
&\quad = (1-x^2)^{s}\sum_{k=0}^\infty b_k^{[\pm s]}\left[\prod_{j=0}^{2s-1} (k+s\pm m-j)+O(1-x)\right](1-x)^{\frac12(-s\pm m)-s+k} \\
%
&\quad = \sum_{k=0}^\infty b_k^{[\pm s]}\left[2^{s}\prod_{j=0}^{2s-1} (k+s\pm m-j)+O(1-x)\right](1-x)^{\frac12|-s\pm m|+k},
\end{align*}
so, combining Lemmas~\ref{lem:angular-commutation}, \ref{lem:angular-ode-L} and Definition~\ref{def:angular-basis}, we deduce
\begin{align*}
s+ m \geq 2s\geq 1\implies \mathfrak B_s^{(1)}= 2^{s}\prod_{j=0}^{2s-1} (s+ m-j), \\ s- m \geq 2s\geq 1 \implies \mathfrak B_s^{(3)}=2^{s}\prod_{j=0}^{2s-1} (s- m-j);
\end{align*}
\item otherwise, if $1\leq s\pm m \leq 2s-1$, then 
\begin{align*}
&(\sin\theta)^{2s}\left(\frac{L^{\pm}_s}{\sin\theta}\right)^{2s} \Xi^{[\pm s]}_{\cN^+} \\
 &\quad =(\sin\theta)^{2s}\left(\frac{L^{\pm}_s}{\sin\theta}\right)^{s\mp m}\left(\frac{L^{\pm}_s}{\sin\theta}\right)^{s\pm m} \Xi^{[\pm s]}_{\cN^+}\\
 &\quad =(\sin\theta)^{2s}\left(\frac{L^{\pm}_s}{\sin\theta}\right)^{s\mp m}\left[(1-x)^{-\frac12(s\pm m)}\sum_{k=0}^\infty b_k^{[\pm s]}\left[\frac{(k+s\pm m)!}{k!}+O(1-x)\right](1-x)^{k}\right]\\
 &\quad =(1+x)^{s}(1-x)^{\frac12|-s\pm m|}\times\\
 &\quad\qquad \times\left(-\frac{d}{dx}-\frac{s\mp m}{2(1+x)}\mp \nu\right)^{s\mp m}\left[\sum_{k=0}^\infty b_k^{[\pm s]}\left[\frac{(k+s\pm m)!}{k!}+O(1-x)\right](1-x)^{k}\right],
\end{align*}
\end{itemize}
so, combining Lemmas~\ref{lem:angular-commutation}, \ref{lem:angular-ode-L} and Definition~\ref{def:angular-basis}, we deduce there are functions $g_\pm$ so that
\begin{gather*}
1\leq s+ m \leq 2s-1\implies \mathfrak B_s^{(1)}= 2^{s}(2s)!\left(b_{s- m}^{[+s]}+g_+(b_0^{[+s]}, \dots, b_{s- m-1}^{[+s]})\right), \\ 1\leq s- m \leq 2s-1 \implies \mathfrak B_s^{(3)}=2^{s}(2s)!\left(b_{s+ m}^{[-s]}+g_-(b_0^{[-s]}, \dots, b_{s+ m-1}^{[-s]})\right).
\end{gather*}

\medskip
\noindent \textit{The $\theta=\pi$ end.} To compute $\mathfrak B_s^{(2)}$ and $\mathfrak B_s^{(4)}$, we can follow the same steps as before. Alternatively, we can  appeal to the symmetry of \eqref{eq:angular-ode} under $\theta\mapsto \pi-\theta$: $L^\pm_s\mapsto -L_s^\mp$  and $\Xi^{[\pm s]}_{\cS^+}(\theta)=\Xi^{[\mp s]}_{\cN^+}(\pi- \theta)$, so the claimed identities follow easily. 

\medskip
\noindent \textit{Non-vanishing of the constants.} Finally, we turn to the proof that $\mathfrak B_s^{(i)}$ are non-zero. In view of the relationships between then, it is enough to show  $\mathfrak B_s^{(1)}$ and $\mathfrak B_s^{(3)}$ are nonzero. Note that, from \eqref{eq:angular-TS-op-rewritten},
\begin{align*}
&\frac{L_s^\pm}{\sin\theta}\left(e^{\mp \nu x}(1-x)^{-\frac12(s\pm m)}(1+x)^{-\frac12(s\mp m)}\right)=0.
\end{align*}
Thus, if we write $\Xi^{[\pm s]} = e^{\mp \nu x}(1-x)^{-\frac12(s\pm m)}(1+x)^{-\frac12(s\mp m)}P^{[\pm s]}(x)$ for some $P^{[\pm s]}\colon (-1,1)\to \mathbb C$, we obtain 
\begin{align*}
\left(\frac{L_s^\pm}{\sin\theta}\right)^{2s}\Xi^{[\pm s]}=(-1)^{2s}e^{\mp \nu x}(1-x)^{-\frac12(s\pm m)}(1+x)^{-\frac12(s\mp m)}\frac{d^{2s}}{dx^{2s}}P_s^{[\pm s]}(x).
\end{align*}

We have established that so far that, if $\Xi^{[\pm s]}$ solves the angular ODE \eqref{eq:angular-ode} and is regular at $\theta=0,\pi$, then it satisfies the identities in the statement for some $\mathfrak B_s^{(1)}$ and $\mathfrak B_s^{(3)}$, respectively. From  the above, it follows that: if $\mathfrak B_s^{(1)}=0$, for regular solutions $\Xi^{[+s]}$, $P^{[+ s]}$ is a polynomial of degree at most $2s-1$; if $\mathfrak B_s^{(3)}=0$, for regular solutions $\Xi^{[-s]}$, then $P^{[- s]}$ must be a polynomial of degree at most $2s-1$. However, the regularity of $\Xi^{[\pm s]}$ forces $P^{[\pm s]}$ to have at least $2s$ zeros (counted with multiplity), so $\mathfrak B_s^{(i)}\neq 0$. Indeed:
\begin{itemize}
\item regularity at $\theta=0$ implies $P^{[\pm s]}$ must vanish to order at least $\frac12|s\pm m|+\frac12(s\pm m)=\max\{s\pm m,0\}$ at $x=1$;
\item regularity at $\theta=\pi$ implies $P^{[\pm s]}$ must vanish to order at least $\frac12|s\mp m|+\frac12(s\mp m)=\max\{s\mp m,0\}$ at $x=-1$;
\end{itemize} 
and thus the polynomial $P^{[\pm s]}$ must have at least $\max\{s\pm m,0\}+\max\{s\mp m,0\}\geq 2s$ roots on the real line.
\end{proof}




\end{document}
